\section{Machine verification and statement alignment}
\label{cert:section}
\begingroup
\raggedright
\setlength{\emergencystretch}{4em}

The submission pins Lean 4.33.1 in \texttt{lean-toolchain}. Its five public roots occupy \src{src/Main.lean}{4}{\texttt{Main.lean}, lines 4--12}. At the level of the stored declaration, the first root has type \texttt{RiemannHypothesis} and value \texttt{rhOfAll}. The second has type \texttt{RiemannZetaZerosOnLine} and the same value. The third has type \texttt{DirichletLNonvanishingRightOfHalf} and value \texttt{allChi}. The fourth has type \texttt{DirichletLZerosOnLine}, with value given by the strip projection of \texttt{claimAll}. The fifth has type \texttt{PrimitiveDirichletLZerosOnLine} and value \texttt{critAll}. The explicit zeta statement excludes negative even integers and $s=1$. The half-plane statement excludes the pair of the principal character and $s=1$. The two Dirichlet critical-line statements are restricted to the open critical strip. The last proposition additionally restricts to primitive characters. The hypotheses are given in \src{src/Statements.lean}{3}{\texttt{Statements.lean}, lines 3--15}. These identifications use the type and proof-value records of the fixed export and the matching source declarations.

\paragraph{Full independent checks and local dependency audit.}
Sutherland reports that all 175,187 compiled modules passed kernel replay. Nanoda and ConLeche accepted the full export, ConLeche compared all five statements with the challenge, and Lean4Lean checked all the compiled modules, including five reruns with a larger step limit \cite{mit}. The record dated 21 November at 09:30 UTC still listed the full-export Lean-kernel and end-to-end Comparator runs as in progress. Its Nanoda log reports a peak resident set of 88,880,436~kB. The local dependency audit traversed constants in declaration types and values together with inductive, constructor, and recursor metadata. The respective root-inclusive closures contain $1{,}478{,}219$, $1{,}478{,}221$, $1{,}478{,}211$, $1{,}478{,}224$, and $1{,}478{,}214$ declarations. The public LUMI table independently gives the first and third counts. Each closure contains precisely the axioms \texttt{Classical.choice}, \texttt{propext}, and \texttt{Quot.sound}. The local graph audit computes dependency reachability from the stored terms. The full proof-checking verdict comes from the public independent runs. See \url{https://math.mit.edu/~drew/grh/} and \url{https://math.mit.edu/~drew/grh/nanoda.log}.

\paragraph{Independent local replay of original declarations.}
The extraction program under \texttt{verification/} checks the SHA-256 of the canonical compressed export, then copies its decompressed bytes from the first record through the original theorem record for \texttt{Deduction.landau\_window}. The resulting prefix has 1,539,032,259 bytes and 140,192 declaration records. Its SHA-256 is \texttt{dfabcd8e32844354b199a7f360c1f11ff}\allowbreak\texttt{e8d7c0bcf67ab6e04c7d36ddbedb08b}. Three original proof values in the prefix mark the scope. The first is \src{src/Formalism/GammaFactorNonvanishing/CritLine.lean}{29}{\texttt{Formalism.VW.rh\_of\_every}}, whose hypothesis is \texttt{UpperTheory.EveryCharacter}. The second is \src{src/Development/AutomorphicArithmetic/DedekindLFunctions/LandauTheorems/poleData.lean}{248}{\texttt{Deduction.poleData}}, giving pole data for the character-family series. The third is \src{src/Development/AutomorphicArithmetic/DedekindLFunctions/LandauTheorems/landauWindow.lean}{83}{\texttt{Deduction.landau\_window}}, which gives $\Re\rho\leq\theta$ for every $\rho\in Z$ from \texttt{PoleData}, $0<\theta<1$, $\delta>0$, and \texttt{WindowHypU}. The prefix contains the exact original records and ends before all five public roots.

The same prefix contains original proof-value records for \src{src/Development/AutomorphicArithmetic/DedekindLFunctions/LandauTheorems/ChCore.lean}{343}{\texttt{Deduction.ch\_core}}, \src{src/Development/AutomorphicArithmetic/DedekindLFunctions/LandauTheorems/ChTop.lean}{37}{\texttt{Deduction.ch\_top}}, and \src{src/Development/AutomorphicArithmetic/PrimitiveCosets/CoreGm.lean}{12}{\texttt{UpperTheory.Comparison}$_2$\texttt{.core\_gm}}. The arithmetic witnesses and their raw export indices are recorded in \path{verification/arithmetic-prefix-receipt.json}. The auxiliary-box density estimate occurs beyond the cut, as do \texttt{UpperTheory.\allowbreak Top.topDesWin} and \texttt{UpperTheory.\allowbreak Gal.Lim.exists\_trace\_limit}.

On native Windows, Nanoda 0.4.17 from the Lean 4.36.0 distribution checked the unmodified prefix and exited successfully. It reported ``Checked 145327 declarations with no errors.'' The larger figure is Nanoda's internal declaration count, while 140,192 counts raw declaration records, each of which may encode an inductive block with multiple constants. The run used a strict allowlist consisting of the three standard axioms above. The success log, configuration, receipts, and reproduction commands are in \path{verification/}. The largest sampled peak working set for this checker process was 1,602,121,728 bytes. The scope of this local kernel verdict is the prefix.

The bundled ConLeche checker independently accepted the same prefix in verified mode with two workers. It reported 140,188 accepted declarations in its own counting convention and completed the 137,728 deferred checks. Its reported elapsed time was 816.7 seconds. The largest sampled peak working set for this process was 2,928,775,168 bytes.

\paragraph{Fresh finite-window interface.}
The file \path{verification/WindowBridge.lean} develops a finite-set induction over the logarithmic interval and then transports it through the exponential map. For an exceptional set $E\subseteq\mathbb N$ with $|E|\le R$, parameters $0<\eta\le C$ and $X>1$, it yields $Y$ such that, with $J=4(2C/\eta)^R$,
\[
 X^{1/J}\le Y,\qquad Y^2\le X,\qquad
 \forall x\in[Y,Y^2]\ \forall k\in E,\quad k\le x^C\Longrightarrow k\le x^\eta.
\]
Here is the corresponding real-variable argument. Put $B=C/\eta$, $U=\log X$, and send each positive exceptional integer $k$ to $a=\log(k)/C$. If there is an exceptional $a\in(U/(2B),U]$, remove it and recurse with upper bound $U/(2B)$. The recursion loses a factor $2B$ and uses one unit of the cardinal bound. If that interval is empty, take $y=U/2$. For each $u\in[y,2y]$, any exceptional $a\le u$ then has $a\le U/(2B)$, so $Ba\le u$. Exponentiating the selected $y$ yields the claim. The case $k=0$ follows directly from positivity of $x^\eta$.

The constant is the source constant since $4(2B)^R=2^{R+2}B^R$. The conclusion agrees with \src{src/Development/AutomorphicArithmetic/AutomorphicVarieties/LayeredIsomorphisms/Window1.lean}{6}{the logarithmic selection lemma} and \src{src/Development/AutomorphicArithmetic/AutomorphicVarieties/LayeredIsomorphisms/Window.lean}{7}{its power-window form}, using the definitions of \src{src/Development/AutomorphicArithmetic/AutomorphicVarieties/PrimeSetBoundDataOfPrimitiveChar.lean}{30}{\texttt{InRange}} and \src{src/Development/AutomorphicArithmetic/AutomorphicVarieties/LayeredIsomorphisms/SecondBound.lean}{46}{\texttt{Good}}. The cardinal bound on $E$ is an explicit hypothesis, and delivering such a bound is an upstream arithmetic obligation. The fresh file was checked by Lean 4.36.0 against Mathlib release \texttt{v4.36.0}. Its only imports are Mathlib modules. Lean's \texttt{\#print axioms} reports precisely \texttt{propext}, \texttt{Classical.choice}, and \texttt{Quot.sound} for each of the two bridge theorems.
\paragraph{Fresh critical-line implication.}
The separate file \path{verification/CriticalLineBridge.lean} checks the quantifiers in Section~\ref{main:end} using Mathlib alone. For a fixed real function it assumes the arithmetic windows at every sufficiently small positive \(\varepsilon\), and assumes the analytic exclusion at each \(0<\theta<1\) for that same function. The Lean proof picks \(\theta\) from the proposed zero's real part, then chooses one \(\varepsilon\) and transports the bound at every point of the resulting windows. Its conclusion excludes every detected point to the right of one half. A separate declaration checks \(M/r^k<c/d^{k+1}\) eventually for \(0<d<r\), \(c>0\), as used to compare a nonzero residue to a Cauchy remainder. All four fresh declarations passed Lean 4.36.0 with the three standard axioms. The analytic exclusion and the arithmetic supply in the implication retain exactly their stated status as hypotheses.
\paragraph{Finite averaging.}
The accompanying \path{verification/TranslateMean.lean} checks two parts
of Proposition~\ref{ar:mean} in Mathlib. For any homomorphism
\(\iota:A\to G\) from a finite group and any additive character
\(\chi\) of \(G\) with values in a field, it computes
\(\sum_{u\in A}\chi(t+\iota(u))\). The sum is
\(|A|\chi(t)\) when the restricted character is trivial, and zero otherwise.
The second proof transfers a uniform absolute error through a normalized
finite average of real numbers. These are formal checks of the
projection and averaging steps. The signs of the arithmetic characters,
the estimates for the three discarded twists, and the geometric input
are justified at their respective places in the text.
\par\endgroup

\subsection{Proof values of the two deformation counts}\label{audit:countpaths}

The proof value of the cohomological window theorem directly uses both
the rank-three Case I count and the joint rank-four Case II and III
count. The accompanying archive \path{verification/zeta-count-trace.zip}
gives the original names and expression IDs for the steps described
in Section~\ref{zetaring:section}. In particular it records the
ternary single-character ring construction, the quartic test detecting
nonfiniteness, and the specialization killing the tower ideal in the
Case II component argument.

The 96 selected declarations lie in the classical RH dependency tree.
All 99 direct references have two witnesses: a constant in the original
compiled value, and a path from the value root to that constant in the
original serialized export. Of these edges, 85 are also immediate
value edges of the saved breadth-first tree. That tree retains one
parent for each declaration. The archive contains 2,935 original
records, including the complete name chains and the type and value
indices for all 95 definitions and theorems. The remaining declaration
is the constructor \(\mathrm{And.intro}\). The verifier checks every
path and can re-extract the original modules from their pinned
archives. The zeta host calculations and the mixed component mean
criterion are paper arguments.

The following pairs are expression IDs in the unmodified export.
The last column counts child-expression steps inside one stored
value. The full sequences, including the two ends of each path,
are in \path{proof-edges.tsv} in the archive.
\begin{center}\small
\begin{tabular}{@{}lrrr@{}}
\toprule
Reference in the stored value & Value root & Constant leaf & Steps\\
\midrule
Window to ternary count                 & 991824206 & 991824173 & 12\\
Window to joint quartic count           & 991824206 & 991824164 & 7\\
Quartic count to conjunction constructor & 952152235 & 1242      & 14\\
Nonfinite test to ambient point         & 962691518 & 962667923 & 57\\
Case II propagation to its last step    & 962437115 & 962426112 & 157\\
\bottomrule
\end{tabular}
\end{center}
